%Old fundamental lemmas about right handed crossings

To recognize whether a crossing is essential and right-handed, it will
be helpful to look just at two different lifts of~$x$.

\begin{lemma}
  Let $a,x \in \pi_1(S)$ be elements that do not share a common power.
  Then $a$ and~$x$ cross essentially iff the lifts $\wt x$ and
  $a \cdot \wt x$ are parallel, and not anti-parallel or crossing.
  Furthermore, if they cross essentially, then $a$ and~$x$ have a
  right-handed crossing iff $a \cdot \wt x$ is on the right-hand side
  of $\wt x$, when viewed along the orientation of~$\wt x$.
  \label{lem:rhs}
\end{lemma}

\DPTtodo[inline]{In Figure~\ref{fig:rhc}, mark and label some
  basepoints on $\wt x$ and $a \cdot \wt x$, and add arrows. Also the
  tildes are not agreeing. Maybe also add the lift $\wt a$ to better
  illustrate N-S dynamics}
  \DMGtodo[inline]{This lemma will be a corollary of the more general arguments on right-handed crossings we have discussed recently. I think the figure will be replaced by the illustration of the more general arguments.}
  \DMGtodo[inline]{Endpoints of different lifts of $\tilde{x}$ can't be anti-parallel, since all elements are hyperbolic (we are dealing with closed surfaces)}
  \begin{figure}
\resizebox{60mm}{!}{\svgfig{righthandedcrossing}}
\caption{Determining a right-handed crossing between $a$ and $x$ by looking at the relative position of two lifts of $x$ in the universal cover.}
\label{fig:rhc}
\end{figure}

\begin{proof}
  This follows easily by consider the North-South dynamics of $a$
  acting on the circle at infinity of~$\wt S$. See Figure~\ref{fig:rhc}.
\end{proof}



\item We take a basepoint $*$ on of the lifts of $x$.
  Let $p \in \pi_1(S,*)$ be sending a point $*'$ in a lift of $C$ to the basepoint $*$, and $q$ an element sending $*'$ to the translate of the basepoint $*$ in another lift of $x$, in such a way that the element $pq$ translates from one lift of $x$ to the other.
    Using the box notation $[p,q,r,s]$, we consider, for each $n$, the boxes associated to
    
    \[
      [p_n,q_n,r_n,s_n] \coloneqq [p,y^{-n}q,yp,y^{-n-1}r]
    \]
    
    See Figure~\ref{fig:recwindow}.
    
           \begin{figure}
\centering{
\resizebox{60mm}{!}{\svgfig{rec_window1}}
\caption{Configuration of lifts and elements of $\pi_1(S)$ used to approximate a fundamental domain of the action on $G(\Sigma)$ for a given curve $C=[y]$.}
\label{fig:recwindow}
}
\end{figure}
    
    This determines, for each $n$, boxes of geodesics whose startpoints are confined in a fixed interval between the startpoints of two lifts of $x$ differing by a translate by $y$, and the endpoints are contained in interval growing as a function of $n$, with endpoints converging to the endpoints of a lift of $[y]$.
    
    Note that, in the RH lift system format $(a,b,c,d)$ it is given by
    \[
      (a_n,b_n,c_n,d_n) \coloneqq (py^nq, pY^{n+1}q,py^{n+1}r,py^{-n}q)
    \]
    
    here we are using $y$ as the element $\pi_1(S,*)$ with red axis in Figure~\ref{fig:recwindow}.
    \item 
    \begin{comment}
    \begin{proposition}
       Let $x$ be the element of $\pi_1(S,*)$ corresponding to a lift of some simple closed curve, $y$ any lift of the closed curve $C$. Let $p,q \in \pi_1(S,*)$ as in Figure~\ref{fig:dreamcatcher}.
       Then,
       
       \[
         f(C)  =  i(\mu_f,C) \lim_{n_1,n_2} \lim_{m_1,m_2} \frac{1}{2} \left(f([py^{n_1+1}rx^{m_1}][py^{-n_2-1}qx^{m_2}])-f([py^{n_1}rx^{m_1}][py^{-n_2}qx^{m_2}]) \right)
       \]
           \end{proposition}
       \end{comment}
       
       \begin{figure}
\centering{
\resizebox{60mm}{!}{\svgfig{dreamcatcher}}
\caption{Configuration of lifts showing that intersection number with $\mu_f$ is defined by a limit of right handed boxes.}
\label{fig:dreamcatcher}
}
\end{figure}
       

       
       

    
    Note that we have
    
    \[
       i(\mu_f,C) = \mu_f((\gamma_-,\gamma_+) \times [x_-,y \cdot x_-))
    \] 
    
    \begin{claim}
If $y^{-1}zy$ and $z$ are parallel then $y$ crosses $z$ to the right.
\label{lem:conj1}
\end{claim}
\begin{proof}
Let $z=ab$, so that $ba=y^{-1}z y$, and $a=y$. Thus, $b=y^{-1}z$
Then, according to Lemma~\ref{lem:righthandedfoundation} Part I, $a$ crosses $b$ to the right, so $y$ crosses $y^{-1}z$ to the right, i.e., $y^{-1}z$ crosses $y^{-1}$ to the right, i.e., $z^{-1}y$ crosses $y$ to the right, and by Lemma~\ref{lem:rhs_prod}, $z^{-1}$ crosses $y$ to the right, so $y$ crosses $z$ to the right
The other option is that $z=aB$, and $y^{-1}z y$, so that $a=y$ and $B=y^{-1}z$, which leads to $Ba$ crossing $aB$ to the right, i.e., $y^{-1}zy$ crossing $z$ to the right, a contradiction.
Since $y^{-1}zy$ and $z$ are conjugate, we can always reduce to one of these two cases, and the proof is concluded. \DMGtodo[inline]{Check this.}
\end{proof}

\begin{claim}
 if $z$ crosses $y^{-1} z y$ to the left, then $z$ is parallel to $y$.
 \label{lem:conj2}
\end{claim}
\begin{proof}
\DMGtodo[inline]{TBW. Should follow by the right handed crossing lemmas.}
\end{proof}
    
    
    %Old proof of $q_1xp_1$ crosses $y$ to the right. It's wrong, but has some interesting ideas that might be useful.
    Suppose first that $q_1 x p_1$ and $y$ are parallel. \DMGtodo[inline]{Need to check the case when they are antiparallel, too, which hopefully will be similar.}
Then $q_1 x p_1 y$ and $y q_1 x p_1$ cross to the left.
If we let $a=q_1x$, $b=p_1 y$, $c=yq_1$ and $d=xp_1$, we have that
$ab$ crosses $cd$ to the left. Conjugating everything by $a$, we get that $ba$ crosses $a^{-1}cda$ to the left. 
Writing things out, we get that $p_1 y q_1 x$ crosses
\[
x^{-1}q_1^{-1}p_1^{-1} (p_1 y q_1 x) p_1 q_1 x.
\]
to the left.
By Lemma~\ref{lem:conj2}, this forces $p_1 y q_1 x$ to be parallel to $p_1q_1x$.
We claim that, in fact, $p_1 y q_1 x$ crosses $p_1q_1x$ to the right.
Indeed, \DMGtodo[inline]{Need to discuss the antiparallel case, too, hopefully similar.} if $p_1 y q_1 x$  was parallel to $p_1q_1x$, by conjugating by $p_1$, it would mean that $y q_1 x p_1$ is parallel to $q_1 x p_1$.
Then by Lemma~\ref{lem:righthandedfoundation} part (II), we would have that 
$y$ is parallel to $(q_1 x p_1)^{-1} y (q_1 x p_1)$.
By then, by Lemma~\ref{lem:conj1}, $q_1 x p_1$ crosses $y$ to the right, which is a contradiction with the hypothesis.

\section{More recent lemmas}

\begin{proposition}
For $C$ a closed curve, we have
\[
i(\mu_f,C)=f(C).
\]
\label{prop:recoverintersection}
\end{proposition}

  Given a representative $y \in \pi_1(S,*)$ of the conjugacy class~$C$, we will first show that we can use right handed boxes to approximate a fundamental domain for $\supp \mu \times \{ \tilde{y} \}$ in $\DG(\Sigma)$.
Second, we'll show that computing the measure of that fundamental domain yields $f([y])$.
Choose a conjugate $y$ in $[y]$ so that $x$ and $y$ are parallel. Thus, $x$ and $y^n$ are parallel for $n>0$ and anti-parallel for $n<0$.
Pick $p_1, q_1 \in \pi_1(S,*)$, both of them crossing $x$ and $y$ to the right.
For each $n$, define the following elements of $\pi_1$: $p_1^n \coloneqq p_1$, $p_2^n \coloneqq p_1 y$, $q_1^n \coloneqq y^n q_1$, and $q_2^n \coloneqq y^{-(n+1)} q_1$.
\begin{proposition}
The datum $[p_1^n,p_2^n,q_1^n,q_2^n]$, as defined above, is a rhs system, i.e.,
\begin{enumerate}
    \item $p_i^nq_j^n$ cross $x$ to the right, for $i,j=1,2$.
    \item $q_2^n x^m p_1^n$ crosses $q_1^n x^m p_2^n$ to the right, for $m \geq 0$. 
\end{enumerate}
\end{proposition}
\begin{proof}
The first part follows by Lemma~\ref{lem:sides}, and the second by Lemma~\ref{lem:diagonals}.
\end{proof}


\begin{lemma}
 $p_1 y^n q_1$ crosses $x$ to the right, for all $n \in \mathbb{Z}$.
 \label{lem:sides}
\end{lemma}
\begin{proof}
For $n=0$, $p_1 q_1$ crosses $x$ to the right by the choice of $p_1$ and $q_1$ and Lemma~\ref{lem:rhs_prod}.
Let $n>0$. Since $x$ and $y^n$ are parallel, by Claim~\ref{lem:conj}, it follows that $x^{-1}y^{-n}x$ crosses $x$ to the right.
Since $p_1x$ crosses $x$ to the right, and $x^{-1}q_1$ crosses $x$ to the right, it follows that $(p_1x)(x^{-1}y^{-n}x)(x^{-1}q_1)=p_1y^{-n}q_1$ crosses $x$ to the right. 
Now we assume $n<0$.
Since $x$ and $y^n$ are antiparallel, then by Lemma~\ref{lem:righthandedfoundation}, part (III), $xy^{n}$ crosses $x$ to the right.
Since $p_1x^{-1}$ crosses $x$ to the right by Lemma~\ref{lem:rhs_prod} and $q_1$ crosses $x$ to the right, it follows that $(p_1x^{-1})(xy^{n})(q_1)=p_1y^{n}q_1$ crosses $x$ to the right. This proves the claim.
\end{proof}

\begin{lemma}
$y^{-(n+1)} q_1 x^m p_1$ crosses $y^n q_1 x^m p_1 y$ to the right, for all $n \geq 0$ and $m>0$.
\label{lem:diagonals}
\end{lemma}
\begin{proof}
Let $a=q_1 x^m p_1$ and $b=y^{n+1}$.
We will prove it by induction.
Let $n=0$. Then, $b=y$.
We want to show that $y^{-1}a$ crosses $ay$ to the right.
We claim that $a$ crosses $y$ to the right.
Then, by Lemma~\ref{lem:righthandedfoundation}, part (I), $y^{-1}a$ crosses $ay$, as we wanted to see.
Suppose that the claim is true for $n>0$, i.e., $c_n \coloneqq y^{-(n+1)}a$ crosses $d_n \coloneqq y^n a y$ to the right.
Note that $c_{n+1} = y^{-1} c_n$ and $d_{n+1}=y d_n$.
Since $a$ crosses $y$ to the right, by Lemma~\ref{lem:rhs_prod}, $d_n$ crosses $y$ to the right, i.e., $y^{-1}$ crosses $d_n$ to the right.
Also, by assumption, $c_n$ crosses $d_n$ to the right.
Thus, by Lemma~\ref{lem:rightleftcomb}, $y^{-1} c_n$ crosses $d_n$ to the right.
Since $a$ crosses $y$ to the right, it follows again by Lemma~\ref{lem:rhs_prod},  that $y^{-1}c_n$ crosses $y$ to the right. Thus, by Lemma~\ref{lem:rightleftcomb}, $y^{-1}c_n$ crosses $y d_n$ to the right.
\end{proof}

\begin{claim}
 $q_1 x^m p_1$ crosses $y$ to the right, for $m>0$.
 \label{lem:cross2}
\end{claim}
\begin{proof}
Without lost of generality, we can assume that $m=1$, since the following argument only uses that $x^m$ is parallel to $y$ for $m>0$, which is equivalent to $x$ being parallel to $y$.
By the choice of $p_1,q_1$, we have that $p_1y$ crosses $y$ to the right, and $Yq_1$ crosses $y$ to the right.
Since $x$ and $y$ are parallel, by Lemma~\ref{lem:conj}, part (iv), $Yxy$ crosses $y$ to the right, and thus by Lemma~\ref{lem:rhs_prod}, $p_1xq_1=(p_1y)(Yxy)(Yq_1)$ crosses $y$ to the right.
\end{proof}